\section{Introduction}
\label{sec:intro}

Multi-task benchmarks and the leaderboards built on them steer which models are adopted, published and funded. A leaderboard gives each model a score on every sub-dataset and then applies an aggregation rule---an average, a median, a count of wins---to produce one ranking. A developer who wants a particular model to be ranked first has more than one lever. One of them is to include data related to some of the sub-datasets in the model's training, which raises the model's scores on those sub-datasets \citep{zhou2023cheater,yang2023rephrased}. How much of this \emph{targeted investment} is needed to reach first place, and how that amount depends on the aggregation rule, has not been quantified on real leaderboards.

We study this question on public data without training any model. We (i) define the investment as a change of the target model's own scores on a chosen set of sub-datasets, with a cap on every score, and measure it by the number of sub-datasets touched ($K$), the total score added ($\Lone$) and the largest single gain ($\Linf$); (ii) compute, \emph{exactly}, the minimum investment that makes a given model strictly first under each of a set of aggregation rules (cardinal rules from the mean to the minimum; ordinal rules from Borda to the Condorcet winner), for the models of two public leaderboards taken as target (all models, except a stratified sample for the rank rules on the larger leaderboard); and (iii) report where the answer is ``impossible'' (no plan within the caps works) as a feasibility share rather than hiding it in the averages.

Because what a developer can achieve on a sub-dataset is not observable, the cap on each score is a modelling assumption, and every statement here is made per cap family: from the permissive ceiling of 100 to the current best score on each sub-dataset, to a fixed number of points above the target's own score, to the gains observed between fine-tuned models and their base models on the Open LLM Leaderboard. We state which findings are robust across these assumptions and which are not.

\paragraph{Contributions.}
(1) A definition of targeted investment in score space with explicit assumptions (single actor, static rivals, no transfer between tasks), and an exact solver for the minimum investment under each rule, validated against closed forms and enumeration (Section~\ref{sec:method}). (2) A measurement of the investment needed and of its feasibility, per rule and per cap, for all targets of HELM MMLU and for the Open LLM Leaderboard (all targets for the score rules, samples for the rank rules) (Section~\ref{sec:experiments}). (3) Robustness checks for normalisation, task granularity, task resampling and evaluation noise, and a statement of when the conclusions fail (Sections~\ref{sec:experiments} and~\ref{sec:limits}). The data, code and one-command reproduction are public (Reproducibility statement).
